% !TeX encoding = UTF-8
\documentclass[variant=guia,cover=institutional,mode=screen]{ua-engineering-v6-article}
% !TeX encoding = UTF-8
\UASetUniversity{Universidad Austral}
\UASetFaculty{Facultad de Ingeniería}
\UASetDepartment{Departamento de Computación}
\UASetCampus{Pilar}
\UASetCareer{Ingeniería Informática}
\UASetCommission{Comisión A}
\UASetProfessor{Equipo docente}
\UASetSemester{Segundo cuatrimestre 2026}
\UASetCourse{Sistemas Distribuidos}
\UASetDate{2026}


\UASetClassNumber{04}
\UASetClassTitle{Consenso distribuido}
\UASetDocKind{Guía de ejercicios}

\title{Clase 04 --- Consenso distribuido}
\author{\UAProfessor}
\date{\UADate}

\begin{document}

\section{Indicaciones generales}
Resuelva cada ejercicio declarando supuestos, modelo de fallas, garantía y trade-off. Cuando corresponda, construya una ejecución verificable y vincule la respuesta con evidencia observable.

\section{Nivel 1: fundamentos y reconocimiento}
\begin{uaexercise}
Definir formalmente el problema de consenso. Indicar claramente entradas, salidas y condiciones del sistema.
\end{uaexercise}

\begin{uaexercise}
Explicar con precisión las propiedades:
Agreement, Validity y Termination.  
Dar un ejemplo concreto donde se viole cada una.
\end{uaexercise}

\begin{uaexercise}
Explicar por qué consenso es necesario en sistemas replicados.
Relacionar con la clase anterior.
\end{uaexercise}

\begin{uaexercise}
Dado un sistema con 3 nodos, construir una ejecución donde dos nodos decidan valores distintos.
¿Qué propiedad se viola?
\end{uaexercise}

\begin{uaexercise}
Definir estado univalente y bivalente.
\end{uaexercise}


\section{Nivel 2: ejecuciones y fallas}
\begin{uaexercise}
Explicar por qué el estado inicial puede ser bivalente.
\end{uaexercise}

\begin{uaexercise}
Explicar el rol del scheduler adversario en la prueba FLP.
\end{uaexercise}

\begin{uaexercise}
Construir una ejecución donde el sistema nunca decide.
\end{uaexercise}

\begin{uaexercise}
Explicar por qué agregar timeouts rompe el modelo de FLP.
\end{uaexercise}

\begin{uaexercise}
Describir detalladamente el proceso de elección de líder en Raft.
\end{uaexercise}


\section{Nivel 3: diseño y comparación}
\begin{uaexercise}
Explicar el rol de los términos (terms).
\end{uaexercise}

\begin{uaexercise}
Describir el protocolo AppendEntries.
\end{uaexercise}

\begin{uaexercise}
Definir formalmente la propiedad de Log Matching.
\end{uaexercise}

\begin{uaexercise}
Explicar por qué el commit requiere mayoría.
\end{uaexercise}

\begin{uaexercise}
Construir una ejecución donde un líder cae y otro lo reemplaza.
\end{uaexercise}


\section{Nivel 4: integración y defensa}
\begin{uaexercise}
Describir los roles de Paxos: proposer, acceptor, learner.
\end{uaexercise}

\begin{uaexercise}
Explicar la fase Prepare.
\end{uaexercise}

\begin{uaexercise}
Explicar la fase Accept.
\end{uaexercise}

\begin{uaexercise}
Definir formalmente qué significa que un valor sea elegido.
\end{uaexercise}

\begin{uaexercise}
Explicar por qué dos valores distintos no pueden ser elegidos.
\end{uaexercise}


\end{document}
